% !TeX program = pdflatex
% Kompilacja: najpierw pdflatex psa.tex (dwukrotnie), potem pdflatex shareholder_agreement.tex (dwukrotnie);
% odesłania do paragrafów Umowy Spółki (\ref{art:...}) są pobierane z psa.aux przez pakiet xr-hyper.
\documentclass[12pt,a4paper]{report}
\usepackage[T1]{fontenc}
\usepackage[utf8]{inputenc}
\usepackage{xr-hyper}
\usepackage{basilisk_i18n}
\usepackage{lmodern}
\usepackage{microtype}
\usepackage{setspace}
\usepackage{parskip}
\usepackage{needspace}
\usepackage{tabularx}
\usepackage[normalem]{ulem}
\externaldocument{psa}

\BusinessPlanCompany{Basilisk Systems}
\BusinessPlanWatermarkText{Basilisk Systems}
\BusinessPlanWatermarkColor{black!6}
\BusinessPlanWatermarkFontSize{38}{46}
\BusinessPlanWatermarkAngle{38}
\BusinessPlanAuthorsText{Basilisk Systems --- projekt umowy wspólników}
\BusinessPlanConfidentialText{Projekt umowy --- tekst czysty}
\BusinessPlanFooterText{Basilisk Systems --- umowa wspólników}
\BusinessPlanVersionText{Wersja 0.9.4-C --- umowa wspólników}

\setstretch{1.18}
\setlist[itemize]{topsep=0.2em,itemsep=0.15em,parsep=0.1em,leftmargin=2.0em}
\setlist[enumerate]{topsep=0.25em,itemsep=0.18em,parsep=0.1em,leftmargin=2.2em}
\setlength{\parindent}{0pt}
\setlength{\parskip}{0.55em}
\setlength{\emergencystretch}{4em}
\renewcommand{\arraystretch}{1.18}

\newcounter{paragraf}
\newcommand{\paragraf}[1]{%
  \refstepcounter{paragraf}%
  \phantomsection%
  \addcontentsline{toc}{section}{\S~\theparagraf\ #1}%
  \Needspace{6\baselineskip}%
  \subsection*{\S~\theparagraf\ #1}%
}
\newenvironment{ContractClause}[1]{%
  \paragraf{#1}%
}{%
  \par
}

\newcommand{\field}[1]{\textbf{[#1]}}
\newcommand{\CompanyFull}{Basilisk Systems prosta spółka akcyjna}
\newcommand{\CompanyShort}{Basilisk Systems P.S.A.}
\newcommand{\KSH}{Kodeks spółek handlowych}
\newcommand{\BoilerplateStrike}[1]{\sout{#1}}
\newcommand{\US}{Umowy Spółki}

\hypersetup{
  pdftitle={Basilisk Systems P.S.A. --- projekt umowy wspólników},
  pdfauthor={Basilisk Systems --- projekt umowy wspólników}
}

\begin{document}

\BusinessPlanTitlePage
  {Basilisk Systems}
  {Wersja 0.9.4-C --- umowa wspólników}
  {Umowa wspólników}
  {Projekt umowy wspólników i umowy wykonawczej mechanizmu zwrotnego zbycia akcji}
  {Dokument towarzyszący umowie prostej spółki akcyjnej}
  {0.9.4-C}

\tableofcontents
\newpage

\chapter*{Projekt umowy wspólników}
\addcontentsline{toc}{chapter}{Projekt umowy wspólników}

\begin{center}
{\Large\bfseries UMOWA WSPÓLNIKÓW}\par
\vspace{0.2cm}
{\large\bfseries wraz z umową wykonawczą mechanizmu zwrotnego zbycia akcji}\par
\vspace{0.25cm}
{\LARGE\bfseries \CompanyFull}\par
\vspace{0.6cm}
\end{center}

zawarta w dniu \field{data} w \field{miejscowość} pomiędzy:

\begin{enumerate}
  \item \field{Założyciel A --- imię, nazwisko, adres, PESEL},
  \item \field{Założyciel B --- imię, nazwisko, adres, PESEL},
  \item \field{Założyciel C --- imię, nazwisko, adres, PESEL},
  \item \field{Założyciel D --- imię, nazwisko, adres, PESEL},
  \item \field{Założyciel E --- imię, nazwisko, adres, PESEL},
\end{enumerate}

zwanymi dalej łącznie „Założycielami”, a każdy z osobna „Założycielem”, oraz --- od chwili przystąpienia zgodnie z §~\ref{sha:strony} --- \CompanyFull{} z siedzibą w Pruszkowie, zwaną dalej „Spółką”.

\textbf{Preambuła.} Założyciele zawiązali Spółkę umową prostej spółki akcyjnej z dnia \field{data} („Umowa Spółki”). Umowa Spółki przewiduje, że szczegółowy tryb wykonania mechanizmu zwrotnego zbycia akcji (§~\ref{art:zwrotne-zbycie} ust.~13 \US), pełnomocnictwa i oferty niezbędne do wykonania obowiązkowego zbycia akcji (§~\ref{art:breach} ust.~3 \US) oraz zakaz konkurencji po ustaniu funkcjonalnego zaangażowania (§~\ref{art:konkurencja} ust.~3 \US) określa odrębna umowa wykonawcza. Niniejsza umowa jest tą umową wykonawczą, a ponadto reguluje wzajemne zobowiązania Założycieli związane z ich zaangażowaniem w Spółkę. Umowa nie zmienia Umowy Spółki i nie zastępuje żadnego jej postanowienia.

\begin{ContractClause}{Definicje i stosunek do Umowy Spółki}
\label{sha:definicje}
\begin{enumerate}
  \item Pojęcia pisane wielką literą, niezdefiniowane odmiennie w niniejszej umowie, mają znaczenie nadane im w Umowie Spółki; dotyczy to w szczególności pojęć: Akcjonariusz Funkcjonalny, Założyciel Pierwotny, Odejście Usprawiedliwione, Odejście Dobrowolne, Odejście Zawinione, Dopuszczalny Nabywca, Niedopuszczalny Nabywca, Kryterium Bezpieczeństwa EU/NATO, Wartość Godziwa, Dzień Wyceny, Niezależny Ekspert, Kluczowa Własność Intelektualna, Przedsięwzięcie Produkcyjne, Kwalifikowana Runda Finansowania, Sprawy Zastrzeżone.
  \item „Dzień Odejścia” oznacza dzień ustania funkcjonalnego zaangażowania Akcjonariusza Funkcjonalnego w Spółkę, ustalony zgodnie z §~\ref{sha:odejscie}.
  \item „Akcje Objęte Mechanizmem” oznaczają akcje Akcjonariusza Funkcjonalnego wskazane w Załączniku nr~2 do Umowy Spółki jako objęte mechanizmem zwrotnego zbycia; „Akcje Niezwolnione” oznaczają Akcje Objęte Mechanizmem, które w Dniu Odejścia nie zostały jeszcze zwolnione z obowiązku zbycia zgodnie z §~\ref{art:zwrotne-zbycie} ust.~2 \US; „Akcje Zwolnione” oznaczają pozostałe Akcje Objęte Mechanizmem.
  \item „Nabywca Wskazany” oznacza Dopuszczalnego Nabywcę wskazanego przez Spółkę zgodnie z §~\ref{sha:nabywca}.
  \item Odesłania do paragrafów Umowy Spółki (§~\ref{art:firma} i następne) odnoszą się do Umowy Spółki w brzmieniu obowiązującym w dniu wykonania danego obowiązku; odesłania bez tego oznaczenia dotyczą niniejszej umowy.
  \item W razie sprzeczności między niniejszą umową a Umową Spółki pierwszeństwo ma Umowa Spółki; postanowienia niniejszej umowy wykłada się tak, aby służyły wykonaniu Umowy Spółki. Niniejsza umowa nie tworzy obowiązków korporacyjnych związanych z akcją innych niż przewidziane w Umowie Spółki i nie ogranicza uprawnień organów Spółki.
\end{enumerate}
\end{ContractClause}

\begin{ContractClause}{Strony, przystąpienie Spółki i nowych akcjonariuszy}
\label{sha:strony}
\begin{enumerate}
  \item Umowa wiąże Założycieli od dnia jej zawarcia. Do chwili przystąpienia Spółki uprawnienia zastrzeżone dla Spółki wykonują pozostali Założyciele działający łącznie, a zobowiązania Założycieli wobec Spółki mają charakter zobowiązań na rzecz osoby trzeciej (art.~393 Kodeksu cywilnego).
  \item Spółka przystępuje do umowy przez złożenie oświadczenia o przystąpieniu według wzoru z Załącznika nr~D niezwłocznie po wpisie Spółki do rejestru przedsiębiorców, nie później niż w terminie 14 dni od wpisu. Założyciele zobowiązują się głosować za uchwałami i podjąć czynności potrzebne do przystąpienia Spółki.
  \item Każda osoba, która nabywa akcje Założyciela Pierwotnego w wykonaniu niniejszej umowy albo obejmuje status Akcjonariusza Funkcjonalnego na podstawie zmiany Umowy Spółki, przystępuje do niniejszej umowy przez złożenie oświadczenia o przystąpieniu według wzoru z Załącznika nr~D przed wpisem do rejestru akcjonariuszy. Założyciele i Spółka mogą uzależnić skuteczność zbycia od takiego przystąpienia w zakresie, w jakim dopuszcza to prawo.
  \item Umowa nie wiąże inwestorów obejmujących akcje w Kwalifikowanej Rundzie Finansowania ani innych akcjonariuszy, którzy nie złożyli oświadczenia o przystąpieniu; stosunki z nimi regulują Umowa Spółki i dokumenty transakcyjne.
\end{enumerate}
\end{ContractClause}

\begin{ContractClause}{Status Akcjonariusza Funkcjonalnego i harmonogram zwalniania akcji}
\label{sha:status}
\begin{enumerate}
  \item Każdy Założyciel wskazany w Załączniku nr~2 do Umowy Spółki jako funkcjonalny potwierdza wpisane tam: liczbę Akcji Objętych Mechanizmem, datę rozpoczęcia oraz rolę i minimalne zaangażowanie. Dane te powtarza Załącznik nr~E do niniejszej umowy; w razie rozbieżności rozstrzyga Załącznik nr~2 do Umowy Spółki.
  \item Akcje Objęte Mechanizmem są zwalniane z obowiązku zbycia zgodnie z §~\ref{art:zwrotne-zbycie} ust.~2 \US: po nieprzerwanym upływie 12 miesięcy od daty rozpoczęcia --- 25\% Akcji Objętych Mechanizmem, a następnie w każdym kolejnym miesiącu --- 1/36 pozostałych 75\%, do pełnego zwolnienia po 48 miesiącach. Liczbę akcji zwolnionych zaokrągla się w dół do pełnej akcji; różnica wynikająca z zaokrągleń jest zwalniana w ostatnim miesiącu.
  \item Zaangażowanie uważa się za nieprzerwane mimo urlopu, choroby albo innej usprawiedliwionej przerwy nieprzekraczającej łącznie \field{90} dni w roku; dłuższa przerwa zawiesza bieg okresu zwalniania na czas jej trwania, chyba że Rada Dyrektorów postanowi inaczej.
  \item Spółka prowadzi ewidencję stanu zwolnienia akcji każdego Akcjonariusza Funkcjonalnego i na jego żądanie wydaje w terminie 14 dni zaświadczenie o liczbie Akcji Zwolnionych i Akcji Niezwolnionych na wskazany dzień.
  \item Zmiana roli albo minimalnego zaangażowania wymaga zgody Rady Dyrektorów i nie przerywa biegu okresu zwalniania, jeżeli zaangażowanie pozostaje stałe, osobiste i istotne w rozumieniu §~\ref{art:zwrotne-zbycie} ust.~1 \US.
\end{enumerate}
\end{ContractClause}

\begin{ContractClause}{Zdarzenie uruchamiające i stwierdzenie rodzaju Odejścia}
\label{sha:odejscie}
\begin{enumerate}
  \item Zdarzeniem uruchamiającym procedurę zwrotnego zbycia jest ustanie funkcjonalnego zaangażowania Akcjonariusza Funkcjonalnego przed pełnym zwolnieniem Akcji Objętych Mechanizmem, w szczególności: rozwiązanie albo wygaśnięcie stosunku prawnego będącego podstawą zaangażowania, złożenie rezygnacji, śmierć, trwała niezdolność do pracy albo uchwała Walnego Zgromadzenia o uzgodnionym odejściu.
  \item Akcjonariusz Funkcjonalny i Spółka zawiadamiają się wzajemnie o zdarzeniu uruchamiającym w terminie 7 dni od jego wystąpienia, w formie dokumentowej, wskazując datę ustania zaangażowania i okoliczności istotne dla kwalifikacji Odejścia.
  \item Rodzaj Odejścia (Usprawiedliwione, Dobrowolne albo Zawinione) stwierdza Rada Dyrektorów uchwałą podjętą w terminie 30 dni od zawiadomienia, po umożliwieniu Akcjonariuszowi Funkcjonalnemu złożenia wyjaśnień; dyrektor będący tym Akcjonariuszem albo pozostający z nim w konflikcie interesów nie uczestniczy w głosowaniu. W przypadkach wymagających uchwały Walnego Zgromadzenia (§~\ref{art:zwrotne-zbycie} ust.~4 lit.~d--e \US) Rada Dyrektorów zwołuje Walne Zgromadzenie w terminie 14 dni od zawiadomienia.
  \item Uchwała określa Dzień Odejścia, rodzaj Odejścia, liczbę Akcji Zwolnionych i Akcji Niezwolnionych na Dzień Odejścia oraz cenę należną za Akcje Niezwolnione. Dniem Odejścia jest dzień ustania zaangażowania, a w razie sporu co do tej daty --- dzień doręczenia pierwszego zawiadomienia z ust.~2.
  \item Jeżeli Rada Dyrektorów nie podejmie uchwały w terminie z ust.~3, Odejście uważa się za Odejście Dobrowolne, chyba że Akcjonariusz Funkcjonalny wykaże okoliczności Odejścia Usprawiedliwionego; Spółka może w każdym czasie stwierdzić Odejście Zawinione, jeżeli okoliczności z §~\ref{art:zwrotne-zbycie} ust.~5 \US{} zostały ujawnione po Dniu Odejścia.
  \item Akcjonariusz Funkcjonalny może w terminie 14 dni od doręczenia uchwały zgłosić sprzeciw co do rodzaju Odejścia lub Dnia Odejścia. Sprzeciw kieruje sprawę do negocjacji i mediacji na zasadach §~\ref{art:final} ust.~4 \US; jeżeli spór dotyczy wyłącznie liczby akcji albo ceny, rozstrzyga go Niezależny Ekspert zgodnie z §~\ref{art:wycena} \US. Sprzeciw nie wstrzymuje obowiązku zbycia Akcji Niezwolnionych po cenie emisyjnej, a jedynie wykonanie prawa nabycia Akcji Zwolnionych.
  \item Od Dnia Odejścia do zakończenia procedury Akcjonariusz Funkcjonalny nie zbywa ani nie obciąża Akcji Objętych Mechanizmem inaczej niż w wykonaniu niniejszej umowy; Spółka może zażądać ujawnienia tego ograniczenia w rejestrze akcjonariuszy zgodnie z §~\ref{art:zwrotne-zbycie} ust.~14 \US.
\end{enumerate}
\end{ContractClause}

\begin{ContractClause}{Wskazanie Dopuszczalnego Nabywcy}
\label{sha:nabywca}
\begin{enumerate}
  \item Nabywcę Wskazanego wyznacza Rada Dyrektorów uchwałą w terminie 60 dni od stwierdzenia rodzaju Odejścia. Nabywcą Wskazanym może być w szczególności: inny Założyciel, osoba wskazana do objęcia statusu Akcjonariusza Funkcjonalnego, uczestnik programu motywacyjnego albo inwestor; Nabywcą Wskazanym nie może być Spółka ani osoba działająca na jej rachunek, z zastrzeżeniem §~\ref{art:zwrotne-zbycie} ust.~12 \US.
  \item Nabywca Wskazany musi być Dopuszczalnym Nabywcą i spełniać Kryterium Bezpieczeństwa EU/NATO zgodnie z §~\ref{art:dopuszczalny-nabywca} \US{} oraz przystąpić do niniejszej umowy w zakresie obowiązków związanych z nabywanymi akcjami.
  \item Przed wskazaniem osoby trzeciej Rada Dyrektorów oferuje nabycie Akcji Niezwolnionych pozostałym Założycielom proporcjonalnie do liczby posiadanych przez nich akcji, z prawem objęcia części niewykorzystanej; Założyciele składają oświadczenia w terminie 14 dni. Postanowienie to nie uchybia §~\ref{art:pierwszenstwo} \US.
  \item Jeżeli Rada Dyrektorów nie wskaże Nabywcy Wskazanego w terminie, Akcje Niezwolnione pozostają objęte obowiązkiem zbycia, a Spółka może wskazać nabywcę w każdym późniejszym czasie; Akcjonariusz Funkcjonalny zachowuje w tym okresie prawa korporacyjne zgodnie z §~\ref{art:zwrotne-zbycie} ust.~3 \US.
  \item Jeżeli Akcje Niezwolnione mają zostać nabyte przez kilku Nabywców Wskazanych, uchwała określa podział akcji między nich; oferta i pełnomocnictwo z §§~\ref{sha:oferta} i~\ref{sha:pelnomocnictwo} obejmują każdego z nich w przypadającej mu części.
\end{enumerate}
\end{ContractClause}

\begin{ContractClause}{Nieodwołalna oferta sprzedaży akcji}
\label{sha:oferta}
\begin{enumerate}
  \item Każdy Akcjonariusz Funkcjonalny składa przy zawarciu niniejszej umowy, według wzoru z Załącznika nr~A, nieodwołalną ofertę sprzedaży Akcji Objętych Mechanizmem na rzecz Nabywcy Wskazanego, ze skutkiem na wypadek Odejścia. Oferta obejmuje: (a)~Akcje Niezwolnione --- po cenie emisyjnej; (b)~Akcje Zwolnione --- w przypadku Odejścia Zawinionego po cenie równej 80\% Wartości Godziwej, nie niższej niż cena emisyjna, a w przypadku Odejścia Dobrowolnego po cenie równej 100\% Wartości Godziwej.
  \item Oferta jest nieodwołalna przez czas trwania mechanizmu zwrotnego zbycia i wygasa w części dotyczącej akcji zwolnionych z obowiązku zbycia, gdy w Dniu Odejścia nie zachodzi podstawa do ich nabycia, a w całości --- z chwilą zwolnienia wszystkich Akcji Objętych Mechanizmem albo ustania obowiązku zbycia z innej przyczyny.
  \item Ofertę przyjmuje Nabywca Wskazany przez oświadczenie złożone Akcjonariuszowi Funkcjonalnemu i Spółce w formie wymaganej przez \KSH{} dla zbycia akcji prostej spółki akcyjnej, w terminie: (a)~30 dni od doręczenia uchwały o wskazaniu --- co do Akcji Niezwolnionych oraz Akcji Zwolnionych przy Odejściu Zawinionym; (b)~6 miesięcy od Dnia Odejścia --- co do Akcji Zwolnionych przy Odejściu Dobrowolnym (§~\ref{art:zwrotne-zbycie} ust.~9 lit.~b \US). Bezskuteczny upływ terminu z lit.~b powoduje wygaśnięcie prawa nabycia Akcji Zwolnionych.
  \item Przyjęcie oferty wymaga wskazania w oświadczeniu liczby nabywanych akcji, ceny albo sposobu jej ustalenia oraz rachunku Nabywcy Wskazanego w rejestrze akcjonariuszy. Umowa sprzedaży dochodzi do skutku z chwilą doręczenia przyjęcia; przeniesienie akcji następuje z chwilą wpisu w rejestrze akcjonariuszy zgodnie z \KSH.
  \item Akcjonariusz Funkcjonalny zobowiązuje się w terminie 7 dni od doręczenia przyjęcia oferty złożyć podmiotowi prowadzącemu rejestr akcjonariuszy dyspozycję wpisu oraz podpisać dokumenty wymagane przez ten podmiot. Spółka współdziała przy dokonaniu wpisu.
\end{enumerate}
\end{ContractClause}

\begin{ContractClause}{Pełnomocnictwo do wykonania obowiązku zbycia}
\label{sha:pelnomocnictwo}
\begin{enumerate}
  \item Na zabezpieczenie wykonania obowiązku zbycia każdy Akcjonariusz Funkcjonalny udziela przy zawarciu niniejszej umowy, według wzoru z Załącznika nr~B, pełnomocnictwa do: (a)~złożenia w jego imieniu oświadczenia o zbyciu Akcji Objętych Mechanizmem na rzecz Nabywcy Wskazanego na warunkach z §~\ref{sha:oferta}; (b)~złożenia dyspozycji wpisu w rejestrze akcjonariuszy i podpisania dokumentów wymaganych przez podmiot prowadzący rejestr; (c)~odbioru ceny na rachunek Akcjonariusza Funkcjonalnego wskazany w pełnomocnictwie.
  \item Pełnomocnictwo jest nieodwołalne i nie wygasa ze śmiercią mocodawcy (art.~101 Kodeksu cywilnego), ponieważ jego nieodwołalność jest uzasadniona treścią stosunku prawnego będącego podstawą jego udzielenia; obejmuje ono wyłącznie czynności wymienione w ust.~1 i może zostać wykorzystane wyłącznie po łącznym spełnieniu warunków z ust.~3.
  \item Pełnomocnik może działać dopiero, gdy: (a)~rodzaj Odejścia został stwierdzony zgodnie z §~\ref{sha:odejscie} i --- w razie sprzeciwu --- spór został rozstrzygnięty; (b)~Nabywca Wskazany przyjął ofertę zgodnie z §~\ref{sha:oferta}; (c)~Akcjonariusz Funkcjonalny nie złożył dyspozycji wpisu w terminie z §~\ref{sha:oferta} ust.~5 mimo pisemnego wezwania z dodatkowym terminem 7 dni; (d)~cena została zapłacona albo złożona do depozytu zgodnie z §~\ref{sha:cena}.
  \item Pełnomocnikiem jest osoba wskazana przez Radę Dyrektorów spośród osób niebędących Nabywcą Wskazanym, członkiem jego organu, jego pełnomocnikiem ani osobą z nim powiązaną, tak aby wykluczyć czynność z samym sobą i reprezentowanie obu stron (art.~108 Kodeksu cywilnego); pełnomocnikiem może być radca prawny, adwokat albo notariusz. Spółka zawiadamia Akcjonariusza Funkcjonalnego o osobie pełnomocnika przed dokonaniem czynności.
  \item Pełnomocnictwo wygasa z chwilą zwolnienia wszystkich Akcji Objętych Mechanizmem, wykonania obowiązku zbycia w całości albo ustania obowiązku zbycia z innej przyczyny; Spółka wydaje na żądanie potwierdzenie wygaśnięcia.
  \item Niezależnie od pełnomocnictwa Spółka i Nabywca Wskazany mogą dochodzić wykonania obowiązku zbycia na drodze sądowej, w tym żądać orzeczenia zastępującego oświadczenie woli (§~\ref{art:zwrotne-zbycie} ust.~15 \US).
\end{enumerate}
\end{ContractClause}

\begin{ContractClause}{Cena, zapłata i przeniesienie akcji}
\label{sha:cena}
\begin{enumerate}
  \item Ceną Akcji Niezwolnionych jest cena emisyjna zapłacona za te akcje przy ich objęciu (0,01 zł za akcję serii A, chyba że z rejestru akcjonariuszy wynika inna cena emisyjna).
  \item Wartość Godziwą Akcji Zwolnionych ustala się zgodnie z §~\ref{art:wycena} \US, przy czym Dniem Wyceny jest Dzień Odejścia. Do czasu ostatecznego ustalenia Wartości Godziwej Nabywca Wskazany może przyjąć ofertę ze wskazaniem ceny „według Wartości Godziwej”; termin zapłaty biegnie wówczas od dnia ostatecznego ustalenia wartości.
  \item Cena jest płatna przelewem na rachunek Akcjonariusza Funkcjonalnego w terminie 30 dni od doręczenia przyjęcia oferty albo od ostatecznego ustalenia Wartości Godziwej, w zależności od tego, co nastąpi później. Jeżeli Akcjonariusz Funkcjonalny nie wskazał rachunku albo odmawia przyjęcia zapłaty, Nabywca Wskazany może złożyć cenę do depozytu sądowego albo notarialnego ze skutkiem zapłaty.
  \item Dyspozycja wpisu w rejestrze akcjonariuszy jest składana nie wcześniej niż po zapłacie albo złożeniu ceny do depozytu; podmiot prowadzący rejestr działa na podstawie dokumentów przewidzianych prawem, a niniejsza umowa nie przesądza jego sposobu działania.
  \item Koszty wyceny ponosi się zgodnie z §~\ref{art:wycena} ust.~6 \US; koszty wpisu w rejestrze ponosi Nabywca Wskazany. Każda strona ponosi własne podatki związane ze zbyciem; Spółka informuje strony o obowiązkach płatnika, jeżeli na niej ciążą.
  \item Rozliczenie z §~\ref{art:zwrotne-zbycie} \US{} nie wyłącza roszczeń Spółki z tytułu Odejścia Zawinionego ani obowiązków trwających po Dniu Odejścia (§~\ref{sha:konkurencja}, §~\ref{sha:ip}, §~\ref{sha:poufnosc}).
\end{enumerate}
\end{ContractClause}

\begin{ContractClause}{Śmierć i zmiana kontroli}
\label{sha:smierc}
\begin{enumerate}
  \item Śmierć Akcjonariusza Funkcjonalnego stanowi Odejście Usprawiedliwione. Akcje Zwolnione podlegają dziedziczeniu na zasadach §~\ref{art:dziedziczenie} \US; Akcje Niezwolnione podlegają zbyciu na rzecz Nabywcy Wskazanego po cenie emisyjnej, a obowiązek ten wykonują spadkobiercy albo pełnomocnik zgodnie z §~\ref{sha:pelnomocnictwo}. Spółka może w takim przypadku przyspieszyć zwolnienie całości albo części akcji zgodnie z §~\ref{art:zwrotne-zbycie} ust.~7 lit.~c \US.
  \item Każdy Założyciel wyraża zgodę na objęcie go ubezpieczeniem na życie na rzecz Spółki zgodnie z §~\ref{art:dziedziczenie} ust.~5 \US{} i zobowiązuje się poddać badaniom oraz złożyć oświadczenia wymagane przez ubezpieczyciela; wysokość sumy ubezpieczenia określa Rada Dyrektorów.
  \item W razie zmiany kontroli nad Spółką stosuje się §~\ref{art:zwrotne-zbycie} ust.~10 \US. Za zmianę kontroli uważa się nabycie przez osobę albo grupę osób działających w porozumieniu ponad 50\% głosów w Spółce albo prawa do powoływania większości dyrektorów.
\end{enumerate}
\end{ContractClause}

\begin{ContractClause}{Zakaz konkurencji po ustaniu zaangażowania}
\label{sha:konkurencja}
\begin{enumerate}
  \item Na podstawie §~\ref{art:konkurencja} ust.~3 \US{} każdy Akcjonariusz Funkcjonalny zobowiązuje się, że przez 12 miesięcy od Dnia Odejścia nie będzie bez pisemnej zgody Spółki: (a)~prowadzić, w jakiejkolwiek formie, działalności konkurencyjnej w zakresie produktów i technologii określonych w §~\ref{art:cel} i §~\ref{art:pkd} \US, w szczególności oprogramowania autonomii, percepcji, koordynacji zespołowej i dowodzenia dla systemów bezzałogowych; (b)~świadczyć pracy ani usług na rzecz podmiotu konkurencyjnego ani uczestniczyć w nim kapitałowo powyżej 2\% (z wyjątkiem spółek publicznych); (c)~nakłaniać pracowników, współpracowników, klientów, partnerów Przedsięwzięć Produkcyjnych ani dostawców Spółki do zakończenia albo ograniczenia współpracy ze Spółką.
  \item Zakaz obowiązuje na terytorium Unii Europejskiej, Europejskiego Obszaru Gospodarczego i państw NATO oraz każdego innego państwa, w którym Spółka w Dniu Odejścia prowadzi sprzedaż albo prace rozwojowe.
  \item Z tytułu zakazu Spółka wypłaca Akcjonariuszowi Funkcjonalnemu odszkodowanie w wysokości \field{kwota albo procent} miesięcznie, jeżeli wymagają tego bezwzględnie obowiązujące przepisy prawa (w szczególności przepisy prawa pracy), a w pozostałych przypadkach --- w wysokości ustalonej w umowie będącej podstawą zaangażowania; brak ustalenia oznacza \field{kwota} miesięcznie. Spółka może w terminie 30 dni od Dnia Odejścia zwolnić Akcjonariusza Funkcjonalnego z zakazu, co zwalnia ją z obowiązku zapłaty odszkodowania za okres po zwolnieniu.
  \item Zakaz nie obejmuje działalności ujawnionej i zatwierdzonej w Załączniku nr~3 do Umowy Spółki, udziału w Przedsięwzięciu Produkcyjnym zatwierdzonego zgodnie z §~\ref{art:przedsiewziecia} ust.~7 \US{} ani działalności naukowej i dydaktycznej niezwiązanej z komercjalizacją technologii konkurencyjnych.
  \item Za każde naruszenie zakazu Akcjonariusz Funkcjonalny zapłaci Spółce karę umowną w wysokości \field{kwota}, a w razie naruszenia trwającego --- dodatkowo \field{kwota} za każdy rozpoczęty miesiąc naruszenia; Spółka może dochodzić odszkodowania przewyższającego karę umowną. Naruszenie zakazu w okresie zaangażowania stanowi Odejście Zawinione zgodnie z §~\ref{art:zwrotne-zbycie} ust.~5 lit.~c \US.
\end{enumerate}
\end{ContractClause}

\begin{ContractClause}{Własność intelektualna}
\label{sha:ip}
\begin{enumerate}
  \item Każdy Założyciel potwierdza, że przeniósł na Spółkę prawa stanowiące jego wkład niepieniężny zgodnie z §~\ref{art:wklady} \US{} i Załącznikiem nr~1 do Umowy Spółki, na podstawie dokumentów wymienionych w Załączniku nr~C do niniejszej umowy, oraz że oświadczenia złożone w §~\ref{art:wlasnosc-intelektualna} \US{} są prawdziwe i kompletne.
  \item Wszelkie utwory, wynalazki, wzory, know-how, dane, projekty i inne rezultaty prac powstałe w okresie zaangażowania Założyciela w związku z działalnością Spółki albo z wykorzystaniem jej zasobów stanowią Kluczową Własność Intelektualną Spółki. Założyciel zobowiązuje się: (a)~zawrzeć ze Spółką, przed rozpoczęciem prac, umowę przenoszącą autorskie prawa majątkowe i prawa do wynalazków na wszystkich polach eksploatacji znanych w chwili zawarcia, z prawem do wykonywania praw zależnych, w zakresie, w jakim prawa te nie przechodzą na Spółkę z mocy prawa; (b)~niezwłocznie ujawniać Spółce powstałe rezultaty i współdziałać przy ich ochronie, także po Dniu Odejścia; (c)~nie wykonywać autorskich praw osobistych w sposób utrudniający korzystanie z utworów przez Spółkę.
  \item Wynagrodzenie za przeniesienie praw, o którym mowa w ust.~2, jest objęte wynagrodzeniem z umowy będącej podstawą zaangażowania, chyba że bezwzględnie obowiązujące przepisy wymagają odrębnego wynagrodzenia; w takim przypadku strony ustalą je w umowie przenoszącej prawa.
  \item Założyciel nie wnosi do prac dla Spółki komponentów objętych prawami osób trzecich ani otwartego oprogramowania o licencjach wymienionych w Załączniku nr~3 do Umowy Spółki jako licencje o podwyższonym ryzyku bez uprzedniej zgody Rady Dyrektorów i ujawnia wykorzystane komponenty zgodnie z §~\ref{art:wlasnosc-intelektualna} \US.
\end{enumerate}
\end{ContractClause}

\begin{ContractClause}{Poufność, bezpieczeństwo i obowiązki informacyjne}
\label{sha:poufnosc}
\begin{enumerate}
  \item Założyciele potwierdzają obowiązki poufności i bezpieczeństwa informacji z §~\ref{art:security} \US, w tym okres 10 lat po zakończeniu posiadania akcji (§~\ref{art:security} ust.~3 \US), oraz obowiązki z §~\ref{art:export} \US. Obowiązki te trwają po Dniu Odejścia niezależnie od rodzaju Odejścia.
  \item Każdy Założyciel aktualizuje ujawnienia z Załącznika nr~3 do Umowy Spółki, w tym udziały i funkcje w Przedsięwzięciach Produkcyjnych oraz aktywności konkurencyjne, w terminie 14 dni od zmiany, w formie dokumentowej skierowanej do Rady Dyrektorów. Aktualizacja nie wymaga zmiany Umowy Spółki; jej zatwierdzenie należy do Rady Dyrektorów zgodnie z §~\ref{art:konkurencja} i §~\ref{art:przedsiewziecia} ust.~7 \US.
  \item Założyciel informuje Spółkę o każdej zmianie ustroju majątkowego małżeńskiego, postępowaniu egzekucyjnym dotyczącym akcji oraz innym zdarzeniu mogącym wpływać na przynależność akcji, zgodnie z §~\ref{art:malzonek} \US.
  \item Za naruszenie obowiązku poufności Założyciel zapłaci Spółce karę umowną w wysokości \field{kwota} za każde naruszenie; Spółka może dochodzić odszkodowania przewyższającego karę umowną.
\end{enumerate}
\end{ContractClause}

\begin{ContractClause}{Współdziałanie Założycieli}
\label{sha:wspoldzialanie}
\begin{enumerate}
  \item Założyciele zobowiązują się współdziałać w dobrej wierze przy przeprowadzeniu Kwalifikowanej Rundy Finansowania zgodnie z §~\ref{art:kwalifikowana-runda} ust.~3 \US, w tym zawrzeć zwyczajowe dokumenty transakcyjne oraz złożyć oświadczenia i zapewnienia dotyczące Spółki w zakresie ich wiedzy, na zasadach uzgodnionych w uchwale kierunkowej.
  \item Założyciele zobowiązują się głosować za uchwałami niezbędnymi do wykonania niniejszej umowy, w szczególności za uchwałami o stwierdzeniu Odejścia Usprawiedliwionego w przypadkach uznaniowych (§~\ref{art:walne-zastrzezone} \US), jeżeli zachodzą przesłanki §~\ref{art:zwrotne-zbycie} ust.~4 \US, oraz nie podejmować działań zmierzających do obejścia mechanizmu zwrotnego zbycia, w szczególności przez przeniesienie akcji na podmiot kontrolowany bez przystąpienia tego podmiotu do niniejszej umowy.
  \item Założyciel zamierzający zbyć akcje na rzecz podmiotu w całości przez siebie kontrolowanego (§~\ref{art:okres-ograniczenia-zbywania} ust.~2 lit.~d \US) zapewnia, że podmiot ten przystąpi do niniejszej umowy, a Założyciel pozostaje odpowiedzialny solidarnie za wykonanie obowiązków z niej wynikających.
  \item Założyciele uzgadniają, że w sprawach objętych niniejszą umową będą dążyć do rozstrzygnięć jednomyślnych, a w razie ich braku --- stosować procedury Umowy Spółki; niniejsza umowa nie ustanawia porozumienia co do sposobu głosowania w innych sprawach.
\end{enumerate}
\end{ContractClause}

\begin{ContractClause}{Czas trwania i zmiana umowy}
\label{sha:czas}
\begin{enumerate}
  \item Umowa obowiązuje od dnia zawarcia do dnia, w którym wszystkie Akcje Objęte Mechanizmem wszystkich Akcjonariuszy Funkcjonalnych zostaną zwolnione z obowiązku zbycia albo zbyte w wykonaniu niniejszej umowy, z tym że obowiązki z §~\ref{sha:konkurencja}, §~\ref{sha:ip} i §~\ref{sha:poufnosc} trwają przez okresy w nich określone, a pełnomocnictwo i oferta --- do chwili ich wygaśnięcia.
  \item Umowa wygasa wobec Założyciela, który przestał być akcjonariuszem Spółki, z zastrzeżeniem obowiązków trwających po Dniu Odejścia.
  \item Zmiana lub rozwiązanie umowy wymaga formy pisemnej pod rygorem nieważności i zgody wszystkich Założycieli będących w danym czasie Akcjonariuszami Funkcjonalnymi oraz Spółki; zmiana nie może prowadzić do sprzeczności z Umową Spółki. Zmiana Umowy Spółki w zakresie mechanizmu zwrotnego zbycia stosuje się do niniejszej umowy bez potrzeby jej zmiany, w granicach ust.~4.
  \item Postanowienia niniejszej umowy, które w wyniku zmiany Umowy Spółki stałyby się z nią sprzeczne, tracą moc w zakresie sprzeczności; strony uzupełnią wówczas umowę zgodnie z §~\ref{sha:koncowe} ust.~2.
\end{enumerate}
\end{ContractClause}

\begin{ContractClause}{Postanowienia końcowe}
\label{sha:koncowe}
\begin{enumerate}
  \item Umowa podlega prawu polskiemu.
  \item Jeżeli postanowienie umowy jest nieważne albo niewykonalne, pozostałe postanowienia zachowują moc, a strony zastąpią wadliwe postanowienie rozwiązaniem prawnie dopuszczalnym, możliwie najbliższym jego celowi; dotyczy to w szczególności pełnomocnictwa, oferty i kar umownych.
  \item Spory związane z umową strony poddają w pierwszej kolejności negocjacjom i mediacji na zasadach §~\ref{art:final} ust.~4--5 \US; w razie braku ugody właściwy jest sąd powszechny właściwy dla siedziby Spółki.
  \item Zawiadomienia składa się w formie dokumentowej na adresy poczty elektronicznej wskazane w Załączniku nr~E, a wobec Spółki --- na adres wpisany do rejestru przedsiębiorców; strony aktualizują adresy w terminie 7 dni od zmiany.
  \item Załączniki nr~A--E stanowią integralną część umowy. Umowę sporządzono w \field{liczba} jednobrzmiących egzemplarzach, po jednym dla każdej ze stron.
\end{enumerate}
\end{ContractClause}

\vspace{1cm}
\begin{center}
\begin{tabularx}{\textwidth}{X X}
\field{Założyciel A --- podpis} & \field{Założyciel B --- podpis} \\
\\[1.3cm]
\field{Założyciel C --- podpis} & \field{Założyciel D --- podpis} \\
\\[1.3cm]
\field{Założyciel E --- podpis} & \field{Spółka --- przystąpienie, data i podpisy zgodnie z reprezentacją} \\
\end{tabularx}
\end{center}

\newpage
\section*{Załącznik nr A — Wzór nieodwołalnej oferty sprzedaży akcji}
\addcontentsline{toc}{section}{Załącznik nr A — Wzór nieodwołalnej oferty sprzedaży akcji}

Ja, niżej podpisany \field{imię, nazwisko, PESEL}, Akcjonariusz Funkcjonalny \CompanyFull{} („Spółka”), działając na podstawie §~\ref{art:zwrotne-zbycie} Umowy Spółki oraz umowy wspólników z dnia \field{data}, składam niniejszym nieodwołalną ofertę sprzedaży przysługujących mi akcji serii A oznaczonych numerami \field{od--do} („Akcje Objęte Mechanizmem”) na rzecz Dopuszczalnego Nabywcy wskazanego przez Spółkę zgodnie z umową wspólników, na następujących warunkach:

\begin{enumerate}
  \item oferta obejmuje Akcje Niezwolnione po cenie emisyjnej oraz --- w przypadku Odejścia Zawinionego albo Odejścia Dobrowolnego --- Akcje Zwolnione po cenie określonej odpowiednio w §~\ref{art:zwrotne-zbycie} ust.~8 lit.~b albo ust.~9 lit.~b Umowy Spółki, ustalonej zgodnie z §~\ref{art:wycena} Umowy Spółki;
  \item liczbę Akcji Niezwolnionych i Akcji Zwolnionych ustala się na Dzień Odejścia zgodnie z umową wspólników;
  \item oferta może zostać przyjęta wyłącznie po stwierdzeniu rodzaju Odejścia, w terminach określonych w umowie wspólników, w formie wymaganej przez \KSH{} dla zbycia akcji;
  \item oferta jest nieodwołalna do chwili zwolnienia wszystkich Akcji Objętych Mechanizmem z obowiązku zbycia albo ustania obowiązku zbycia z innej przyczyny;
  \item cena jest płatna na rachunek bankowy nr \field{numer} w terminie określonym w umowie wspólników.
\end{enumerate}

\field{miejscowość, data} \hfill \field{podpis}

\newpage
\section*{Załącznik nr B — Wzór pełnomocnictwa}
\addcontentsline{toc}{section}{Załącznik nr B — Wzór pełnomocnictwa}

Ja, niżej podpisany \field{imię, nazwisko, PESEL}, Akcjonariusz Funkcjonalny \CompanyFull{} („Spółka”), udzielam niniejszym osobie wskazanej przez Radę Dyrektorów Spółki zgodnie z §~\ref{sha:pelnomocnictwo} ust.~4 umowy wspólników z dnia \field{data} pełnomocnictwa do:

\begin{enumerate}
  \item złożenia w moim imieniu oświadczenia o zbyciu Akcji Objętych Mechanizmem, oznaczonych numerami \field{od--do}, na rzecz Dopuszczalnego Nabywcy wskazanego przez Spółkę, na warunkach określonych w mojej nieodwołalnej ofercie z dnia \field{data} i w umowie wspólników;
  \item złożenia dyspozycji wpisu zbycia w rejestrze akcjonariuszy oraz podpisania dokumentów wymaganych przez podmiot prowadzący rejestr;
  \item odbioru ceny sprzedaży na rachunek bankowy nr \field{numer}.
\end{enumerate}

Pełnomocnictwo jest nieodwołalne i nie wygasa z chwilą mojej śmierci (art.~101 Kodeksu cywilnego), a jego nieodwołalność jest uzasadniona treścią stosunku prawnego wynikającego z Umowy Spółki i umowy wspólników. Pełnomocnik może działać wyłącznie po łącznym spełnieniu warunków określonych w §~\ref{sha:pelnomocnictwo} ust.~3 umowy wspólników i nie może być stroną czynności ani osobą powiązaną z nabywcą. Pełnomocnictwo wygasa z chwilą zwolnienia wszystkich Akcji Objętych Mechanizmem z obowiązku zbycia albo wykonania obowiązku zbycia.

\field{miejscowość, data} \hfill \field{podpis notarialnie poświadczony, jeżeli wymaga tego podmiot prowadzący rejestr}

\newpage
\section*{Załącznik nr C — Wykaz dokumentów przeniesienia własności intelektualnej}
\addcontentsline{toc}{section}{Załącznik nr C — Wykaz dokumentów przeniesienia własności intelektualnej}

\begin{longtable}{@{}p{1.3cm} p{4.2cm} p{4.0cm} p{4.6cm}@{}}
\toprule
\textbf{Zał.} & \textbf{Przedmiot wkładu} & \textbf{Dokument przenoszący} & \textbf{Data i strony} \\
\midrule
\endfirsthead
\toprule
\textbf{Zał.} & \textbf{Przedmiot wkładu} & \textbf{Dokument przenoszący} & \textbf{Data i strony} \\
\midrule
\endhead
A & \field{opis zgodny z Załącznikiem nr~1 do Umowy Spółki} & \field{umowa przeniesienia / cesja / protokół wydania} & \field{data} \\
B & \field{opis} & \field{dokument} & \field{data} \\
C & \field{opis} & \field{dokument} & \field{data} \\
D & \field{opis} & \field{dokument} & \field{data} \\
E & \field{opis} & \field{dokument} & \field{data} \\
\bottomrule
\end{longtable}

\newpage
\section*{Załącznik nr D — Wzór oświadczenia o przystąpieniu}
\addcontentsline{toc}{section}{Załącznik nr D — Wzór oświadczenia o przystąpieniu}

\field{Spółka / nabywca akcji / nowy Akcjonariusz Funkcjonalny --- oznaczenie}, oświadcza, że zapoznał się z umową wspólników \CompanyFull{} z dnia \field{data} i przystępuje do niej z dniem \field{data} jako \field{Spółka / nabywca akcji nr od--do / Akcjonariusz Funkcjonalny}, przyjmując wszystkie prawa i obowiązki wynikające z tej umowy w zakresie \field{całości / nabywanych akcji}, w tym --- jeżeli dotyczy --- składając ofertę według Załącznika nr~A i udzielając pełnomocnictwa według Załącznika nr~B.

\field{miejscowość, data} \hfill \field{podpis(y) zgodnie z reprezentacją}

\newpage
\section*{Załącznik nr E — Akcjonariusze Funkcjonalni i dane do zawiadomień}
\addcontentsline{toc}{section}{Załącznik nr E — Akcjonariusze Funkcjonalni i dane do zawiadomień}

Dane powtarzają Załącznik nr~2 do Umowy Spółki; w razie rozbieżności rozstrzyga Załącznik nr~2 do Umowy Spółki.

\begin{longtable}{@{}p{1.3cm} p{2.2cm} p{2.6cm} p{2.4cm} p{2.6cm} p{3.0cm}@{}}
\toprule
\textbf{Zał.} & \textbf{Funkcjonalny?} & \textbf{Akcje objęte (nr od--do)} & \textbf{Data rozpoczęcia} & \textbf{Rola} & \textbf{E-mail do zawiadomień} \\
\midrule
\endfirsthead
\toprule
\textbf{Zał.} & \textbf{Funkcjonalny?} & \textbf{Akcje objęte (nr od--do)} & \textbf{Data rozpoczęcia} & \textbf{Rola} & \textbf{E-mail do zawiadomień} \\
\midrule
\endhead
A & \field{TAK/NIE} & \field{od--do} & \field{data} & \field{rola} & \field{e-mail} \\
B & \field{TAK/NIE} & \field{od--do} & \field{data} & \field{rola} & \field{e-mail} \\
C & \field{TAK/NIE} & \field{od--do} & \field{data} & \field{rola} & \field{e-mail} \\
D & \field{TAK/NIE} & \field{od--do} & \field{data} & \field{rola} & \field{e-mail} \\
E & \field{TAK/NIE} & \field{od--do} & \field{data} & \field{rola} & \field{e-mail} \\
\bottomrule
\end{longtable}

\end{document}
